\section{第41届复赛第2题：台球弹性碰撞、滑动摩擦与30度经验规则}
\subsection{题干}
质量、半径相同的均匀实心台球 $Q$、$T$ 发生完全弹性碰撞，球间无摩擦。碰撞前 $Q$ 球质心速度 $\boldsymbol v_0$ 沿 $y$ 方向，角速度为 $\boldsymbol\omega_0$。碰撞后，$T$ 球以与 $y$ 轴夹角 $\varphi$ 飞出，$Q$ 球以夹角 $\theta$ 飞出；随后 $Q$ 球在桌面摩擦下过渡至纯滚。题目依次要求求碰撞后的速度和角速度、接触点速度、纯滚所需时间与轨迹、达到纯滚后的方向；又要求确定使开球立即纯滚的击球高度，并验证30度经验规则。
\begin{figure}[htbp]\centering\begin{tikzpicture}[>=Stealth,scale=.88]
  \draw[->,thick] (-.4,0)--(5.2,0) node[right] {$x$};
  \draw[->,thick] (1.7,-.45)--(1.7,4.3) node[above] {$y$};
  \draw[thick] (1.7,0) circle(.63); \node[below right] at (2.12,-.12) {$Q$球};
  \draw[thick] (.62,1.15) circle(.63); \node[left] at (-.18,1.02) {$T$球};
  \fill (1.7,0) circle(1.2pt); \fill (.62,1.15) circle(1.2pt);
  \draw[->,very thick] (1.7,0)--(3.55,2.25) node[above right] {$\boldsymbol v_Q$};
  \draw[->,very thick] (.62,1.15)--(-.55,2.65) node[above left] {$\boldsymbol v_T$};
  \draw (1.7,1.1) arc[start angle=90,end angle=50,radius=1.1]; \node at (2.25,1.34) {$\theta$};
  \draw (1.7,1.95) arc[start angle=90,end angle=128,radius=1.95]; \node at (.98,2.52) {$\varphi$};
  \node[below] at (1.7,-.82) {原题图2a：两球碰撞后的飞行方向};
\end{tikzpicture}
\caption{第41届复赛第2题原题图2a的TikZ重绘。}\end{figure}
\begin{figure}[htbp]\centering\begin{tikzpicture}[>=Stealth,scale=.88]
  \draw[->,thick] (-.4,0)--(5.1,0) node[right] {$x$};
  \draw[->,thick] (1.7,-.4)--(1.7,4.3) node[above] {$y$};
  \draw[thick] (1.7,0) circle(.63); \node[below right] at (2.12,-.12) {$Q$球};
  \draw[thick] (.62,1.15) circle(.63); \node[left] at (-.18,1.02) {$T$球};
  \draw[->,very thick] (.62,1.15)--(-.55,2.65) node[above left] {$\boldsymbol v_T$};
  \draw (1.7,1.95) arc[start angle=90,end angle=128,radius=1.95]; \node at (.98,2.52) {$\varphi$};
  \draw[densely dashed] (3.75,.25)--(3.75,3.8);
  \draw[->,very thick] (1.7,0)--(3.75,3.55) node[above right] {$\boldsymbol v_{Qf}$};
  \draw (3.75,2.35) arc[start angle=90,end angle=60,radius=1.2]; \node[right] at (3.88,2.75) {$\theta_f$};
  \node[below] at (1.9,-.82) {原题图2b：$Q$球达到纯滚后的最终方向};
\end{tikzpicture}
\caption{第41届复赛第2题原题图2b的TikZ重绘。}\end{figure}
\begin{figure}[htbp]\centering\begin{tikzpicture}[>=Stealth,scale=.92]
  \fill[gray!20] (-2.4,-1.25) rectangle (2.4,-.95);
  \foreach \x in {-2.25,-2.0,...,2.25} \draw[gray] (\x,-1.25)--++(.18,.3);
  \draw[thick] (0,-.32) circle(1.02);
  \fill (0,-.32) circle(1.2pt); \node[right] at (.1,-.15) {$Q$球};
  \draw[->,very thick] (1.02,-.32)--(2.2,-.32) node[right] {$\boldsymbol v_0$};
  \draw[->,very thick] (-2.0,.5)--(-.85,.5) node[midway,above] {$\boldsymbol M$};
  \draw[densely dashed] (-.85,.5)--(-.85,-1.02);
  \draw[<->] (-1.45,.5)--(-1.45,-1.02) node[midway,left] {$h$};
  \draw[->,very thick] (-.45,.92) arc[start angle=125,end angle=15,radius=.7] node[midway,above] {$\boldsymbol\omega_0$};
  \node[below] at (0,-1.55) {原题图2c：球杆冲量与开球自旋};
\end{tikzpicture}
\caption{第41届复赛第2题原题图2c的TikZ重绘。}\end{figure}
\subsection{参考解答整理}
取碰撞连心线为 $y'$ 方向。由动量守恒和完全弹性条件，碰撞后的两球飞行方向满足
\[\boxed{\theta+\varphi=90^\circ},\qquad \boxed{\boldsymbol\omega=\boldsymbol\omega_0},\]
且 $Q$ 球质心速度分量为 $v_{Qx}=v_0\sin\varphi\cos\varphi$、$v_{Qy}=v_0\sin^2\varphi$。答案中的质心系速度几何关系如图所示。
\begin{figure}[htbp]\centering\begin{tikzpicture}[>=Stealth,scale=.9]
  \coordinate (O) at (0,0);
  \draw[->,thick] (O)--(0,3.1) node[above] {$y'$};
  \draw[->,thick] (O)--(3.8,0) node[right] {$x'$};
  \draw[->,very thick] (O)--(-1.8,1.1) node[above left] {$\boldsymbol v_T$};
  \draw[->,very thick] (O)--(1.45,2.35) node[above right] {$\boldsymbol v_{Q'}$};
  \draw[->,very thick] (O)--(-1.35,-.85) node[below left] {$\boldsymbol v_Q$};
  \draw[->,very thick] (O)--(.82,.48) node[right] {$\boldsymbol v_{CM}$};
  \draw (0,1.12) arc[start angle=90,end angle=58,radius=1.12]; \node at (.45,1.32) {$\varphi$};
  \draw (0,-.92) arc[start angle=180,end angle=212,radius=.92]; \node at (-.68,-.55) {$\theta$};
  \node[below] at (.5,-1.55) {答案图2b：碰撞速度的质心系几何关系};
\end{tikzpicture}
\caption{第41届复赛第2题答案图2b的TikZ重绘。}\end{figure}
碰撞后接触点的初速度为
\[\boldsymbol v_{C0}=(v_0\sin\varphi\cos\varphi-R\omega_{0y})\boldsymbol i+(v_0\sin^2\varphi+R\omega_{0x})\boldsymbol j,\]
故达到纯滚的时间为 $\Delta t=2|\boldsymbol v_{C0}|/(7\mu g)$。在滑动阶段，摩擦力方向与接触点速度相反，质心轨迹为抛物线；达到纯滚后摩擦力消失，质心作匀速直线运动。

开球冲量满足 $M=mv_0$、$(h-R)M=I\omega_0$。对均匀实心球 $I=2mR^2/5$，使开球即纯滚的击球高度为
\[\boxed{h=\frac{7R}{5}}.\]
在此条件下，答案给出
\[\tan\theta_f=\frac{\sin\varphi\cos\varphi}{\sin^2\varphi+\frac25}=\frac{5\sin2\varphi}{9-5\cos2\varphi}.\]
极大值出现于 $\cos2\varphi_m=5/9$，即 $\varphi_m=28.1^\circ$；相应 $\theta_m=33.7^\circ$。题目表格的五个结果为：
\begin{center}
\begin{tabular}{c|cccccc}\hline
$\varphi$ & $\varphi_m$ & $\varphi_m+5^\circ$ & $\varphi_m+10^\circ$ & $\varphi_m-5^\circ$ & $\varphi_m-10^\circ$\\\hline
$\theta_f$ & $33.7^\circ$ & $33.3^\circ$ & $32.0^\circ$ & $33.1^\circ$ & $30.7^\circ$\\\hline
\end{tabular}
\end{center}
\subsection{转写依据}
题干依据 \texttt{CPHO41-Q02\_transcription\_v1.md}；公式和数值依据答案第3至7页及 \texttt{CPHO41-Q02\_formula\_check\_v1.csv}。图2a、图2b、图2c和答案图2b均由原页重绘；没有嵌入原始页面截图。
